% ---- pg-455 （书页 443） ----

出的.

\noindent 10.10\quad 在一项有关贫穷、犯罪和威慑的研究中, 帕克和史密斯(Parker and Smith, 见文献[10])报告了在 1970 年和 1973 年不同州的犯罪统计数据. 部分样本相关系数矩阵为
\[
\mathbf{R}=\left[\begin{array}{c:c}\mathbf{R}_{11}&\mathbf{R}_{12}\\ \hdashline \mathbf{R}_{21}&\mathbf{R}_{22}\end{array}\right]=\left[\begin{array}{cc:cc}1.0&0.615&-0.111&-0.266\\0.615&1.0&-0.195&-0.085\\ \hdashline -0.111&-0.195&1.0&-0.269\\-0.266&-0.085&-0.269&1.0\end{array}\right]
\]
变量为

$X_1^{(1)}=1973$ 年非一级杀人犯

$X_2^{(1)}=1973$ 年一级杀人犯(杀人者为家庭成员和熟人)

$X_1^{(2)}=1970$ 年处罚的严厉性(服刑月份的中值)

$X_2^{(2)}=1970$ 年处罚的确定性(监狱接纳人数除以杀人犯总数)

\noindent (a)\quad 求样本典型相关系数.

\noindent (b)\quad 确定第一典型对 $\hat{U}_1,\hat{V}_1$ 并解释这些量.

\noindent 10.11\quad 例 8.5 中列出了从五支股票的 $n=103$ 连续的周回报率得到相关系数矩阵. 试作出银行回报率 $\mathbf{X}^{(1)}=[X_1^{(1)},X_2^{(1)},X_3^{(1)}]'$ 和石油公司回报率 $\mathbf{X}^{(2)}=[X_1^{(2)},X_2^{(2)}]'$ 的典型相关分析.

\noindent 10.12\quad 对 $n=70$ 个家庭的随机样本进行调查, 用以确定某“人口统计”变量和某“消费”变量的联系.

令:

判别组 $\left\{\begin{array}{l}X_1^{(1)}=\text{每年去餐馆就餐的频率}\\X_2^{(1)}=\text{每年出外看电影的频率}\end{array}\right.$

预报组 $\left\{\begin{array}{l}X_1^{(2)}=\text{户主的年龄}\\X_2^{(2)}=\text{家庭年收入}\\X_3^{(2)}=\text{户主受教育程度}\end{array}\right.$

假设对以上变量的 70 个观测值给出样本相关系数矩阵
\[
\mathbf{R}=\left[\begin{array}{c:c}\mathbf{R}_{11}&\mathbf{R}_{12}\\ \hdashline \mathbf{R}_{21}&\mathbf{R}_{22}\end{array}\right]=\left[\begin{array}{cc:ccc}
1.0& & & & \\
0.80&1.0& & & \\ \hdashline
0.26&0.33&1.0& & \\
0.67&0.59&0.37&1.0& \\
0.34&0.34&0.21&0.35&1.0
\end{array}\right]
\]

\noindent (a)\quad 确定样本典型相关系数, 并且在显著性水平 $\alpha=0.05$ 下, 检验假设 $H_0:\boldsymbol{\Sigma}_{12}=\mathbf{0}$(或, 等价地, $\boldsymbol{\rho}_{12}=\mathbf{0}$). 若拒绝 $H_0$, 检验第一典型相关系数的显著性($\alpha=0.05$).

\noindent (b)\quad 用标准化变量, 构造与“显著的”典型相关系数相对应的典型变量.

\noindent (c)\quad 利用(a)及(b)的结果, 将典型变量系数(对“显著的”典型相关系数)和典型变量与其成分变量的样本相关系数列表.

\noindent (d)\quad 就(c)给出的信息, 阐述这些典型变量.

\noindent (e)\quad 人口统计变量对消费变量是否有影响? 消费变量比人口统计变量提供的信息要多吗?
